%%%%%%%%%%%%%%%%%%%%%%%%%%%%%%%%%%%%%%%%%%%%%%%%%%%%%%%%%%%%%%%%%%%%%%%%%%%%%%%
% SECTION: DISCUSSION (IEEE TSE Format) — Consolidated
%%%%%%%%%%%%%%%%%%%%%%%%%%%%%%%%%%%%%%%%%%%%%%%%%%%%%%%%%%%%%%%%%%%%%%%%%%%%%%%

This section synthesizes the implications of our empirical findings, analyzes the
architectural factors driving \filopriori{}'s effectiveness, and outlines practical
considerations for its deployment in CI/CD environments.

\subsection{Structural Dominance: Why Graph Modeling Matters and Semantic Features Do Not}

The ablation studies on both datasets provide converging evidence that
modeling co-failure relationships between test cases through GATv2 is the
most impactful design decision for TCP. On the industrial dataset, removing
the GATv2 graph attention mechanism causes a 17.0\% performance decrease
($p < 0.001$), and the per-edge-type ablation
(Table~\ref{tab:edge_ablation}) confirms that co-failure edges are the
decisive graph component: a co-failure-only graph achieves APFD = 0.763,
matching or exceeding the full 5-type configuration. On the RTPTorrent
dataset, the GNN architecture maintains strong performance using only
co-failure edges, while the DNN ensemble provides a complementary
contribution ($-$2.6\%, $p < 0.001$). We attribute the effectiveness of
co-failure graph modeling to several factors.

First, graph-based modeling encodes relationships that flat feature vectors
cannot capture: tests that co-fail often share underlying dependencies on the
same code modules, and the GAT learns to propagate failure signals through
these connections. Second, GATv2~\cite{brody2022attentive} computes dynamic
attention that depends on both query and key nodes, allowing the model to
selectively attend to the most relevant neighbors for each test case, adapting
to different failure patterns.

Third, the per-edge-type ablation (Table~\ref{tab:edge_ablation}) reveals
that co-failure edges alone provide sufficient structural signal for the
GATv2 attention mechanism. With all other parameters held constant,
co-failure-only graphs slightly outperform 5-type graphs (0.763 vs.\ 0.754,
$p = 0.002$), suggesting that supplementary edge types may introduce noise
that dilutes the co-failure signal. This refines the interpretation of the
+10.0\% attributed to the ``Enriched Multi-Edge Graph'' in the full-system
ablation (Table~\ref{tab:ablation}): that improvement was primarily driven
by graph construction parameter optimization rather than edge type diversity. This finding strengthens the
deployment case: co-failure edges are universally available in any CI/CD
system, requiring no additional metadata. On RTPTorrent, where only
co-failure edges are available, the GNN still outperforms flat-feature
baselines, confirming that single-edge-type graphs capture meaningful
structural signal.

The finding that simpler GNN architectures (1 layer, 2 heads) outperform deeper
ones (2--3 layers, 4--8 heads) suggests that the test relationship graph can be
effectively modeled with shallow message passing. This is consistent with the
observation that most test case relationships are captured within one hop on
the graph, and deeper networks may oversmooth node representations.

A related finding is that generic, pre-trained sentence-level
embeddings are subsumed by structural execution patterns for test case
prioritization, even when rich textual metadata is available. On the
industrial dataset, removing the semantic stream (Sentence-BERT) reduces
APFD by only 1.9\% ($p = 0.309$, not significant), while removing the
structural stream causes a 5.3\% drop ($p = 0.007$). On RTPTorrent,
removing the semantic stream changes APFD by exactly 0.0\% on 19 of 20
projects.

This result is noteworthy because it qualifies the assumption in TCP
literature that semantic similarity between tests and code changes is a
strong predictor of failure~\cite{hernandes2025method}. Our evidence
suggests that when historical execution features (failure rates, recency,
trends) and graph-based co-failure propagation are available, they capture
the discriminative signal more effectively than domain-agnostic textual
embeddings. More precisely, how tests have behaved is more
predictive than what generic sentence encoders capture about test
descriptions. This finding applies specifically to pre-trained,
domain-agnostic models (Sentence-BERT with \texttt{all-mpnet-base-v2});
code-aware language models (e.g., CodeBERT~\cite{feng2020codebert},
StarCoder), AST-based representations, or embeddings fine-tuned on
software engineering corpora may extract deeper semantic signals that
generic encoders miss. Our methodology (particularly the Random-Fixed
control experiment in Table~\ref{tab:semantic_control}) provides a
principled way to distinguish genuine semantic contributions from identity
proxy effects in future evaluations of such models.

We identify three explanations for this finding, supported by both
ablation evidence and a dedicated control experiment
(Table~\ref{tab:semantic_control}). First, structural features such as
\texttt{recent\_failure\_rate} and \texttt{consecutive\_failures} directly
encode the temporal patterns that predict future failures, while semantic
similarity between test descriptions captures only indirect functional
relationships. Second, the GATv2 attention mechanism over the co-failure
graph already propagates failure signals between functionally related
tests, partially subsuming the role that semantic similarity would play.
Third, even the ``rich'' industrial descriptions may lack the specificity
needed to predict failures. For example, a test described as ``verify OTA
download'' may fail for reasons unrelated to its description (e.g.,
infrastructure issues, memory leaks).

\textbf{The Identity Proxy Effect.}
The control experiment in Section~\ref{sec:rq1_rtptorrent} provides
the strongest evidence for this finding. FailRank-BB achieves competitive
performance on RTPTorrent (APFD = 0.822) using only SBERT embeddings of
test names, which might appear to demonstrate the value of semantic
features. However, replacing SBERT embeddings with random but
consistent vectors (Random-Fixed) yields an even higher APFD (0.856),
while destroying identity consistency (Random-Shuffled) collapses
performance to the random baseline (0.514). This result demonstrates that
FailRank-BB's Logistic Regression classifier is not learning from semantic
content. It is using the embedding as a unique fingerprint for each
test case, effectively learning a per-test failure rate lookup table from
the training data.

This identity proxy mechanism explains the paradoxical cross-dataset
behavior of semantic approaches:
\begin{itemize}
    \item On \textbf{RTPTorrent}, each project has a bounded set of test
    cases that recur across builds. The SBERT embedding provides a
    consistent identifier per test, allowing the classifier to memorize
    per-test failure probabilities. Random-Fixed vectors perform even
    better because they avoid the confounding effect of SBERT placing
    semantically similar (but behaviorally distinct) tests close together
    in the embedding space.
    \item On the \textbf{industrial dataset}, FailRank-BB is the weakest
    baseline (APFD = 0.595). The larger and more dynamic test suite
    (2,347 unique test cases), combined with more complex failure patterns
    and severe class imbalance (37:1), makes the identity-based lookup
    less effective: the Logistic Regression cannot reliably separate
    failure-prone tests in the high-dimensional embedding space when the
    training signal is sparse and noisy.
\end{itemize}

This analysis also explains why \filopriori{}'s semantic stream provides
no significant improvement ($p = 0.309$): the 19 structural features
(including \texttt{failure\_rate}, \texttt{recent\_failure\_rate}, and
\texttt{consecutive\_failures}) already capture the per-test failure
signal explicitly, making the implicit identity-based signal
from embeddings redundant. The GATv2 graph further propagates this
signal between related tests, leaving no residual information for the
semantic stream to contribute. We confirmed this hypothesis by applying
the Random-Fixed control experiment to \filopriori{} itself
(Table~\ref{tab:filopriori_randomfixed}): replacing SBERT embeddings
with hash-seeded random vectors yields APFD = 0.756, statistically
indistinguishable from the SBERT version (0.756, $p = 0.965$). This
closes the identity proxy analysis: the semantic stream acts as an
identity proxy not only in FailRank-BB but also in \filopriori{}'s
own architecture.

The cross-attention fusion module shows a related pattern: removing it
slightly improves the base architecture's APFD (from 0.6397 to 0.6466,
$+$1.1\%, $p = 0.155$, not significant). This suggests that when the
semantic stream carries limited discriminative signal, the cross-attention
mechanism has little meaningful content to fuse and may introduce noise.
We retain cross-attention as an architectural option whose contribution
may emerge in domains with stronger textual-behavioral correlations
(e.g., test descriptions that directly reference code modules under
change), but acknowledge that in our evaluation it does not provide a
statistically significant benefit.

\textbf{Practical implication:} Practitioners deploying \filopriori{}
should prioritize collecting high-quality execution history over investing
in detailed test descriptions when using generic sentence encoders.
The structural features and graph-based modeling provide the dominant
signal. Semantic information extracted via domain-agnostic models like
Sentence-BERT offers marginal improvements even in the most favorable
setting, and when it appears to help (as in FailRank-BB on RTPTorrent),
the actual mechanism is identity memorization rather than semantic
understanding. This does not imply that all NLP-based approaches are
ineffective for TCP; rather, it establishes that the bar for semantic
models is high when rich execution history is available. Future
extensions using code-aware language models fine-tuned on software
engineering corpora might extract deeper semantic signals. Our control
experiment (Table~\ref{tab:semantic_control}) provides a principled
methodology for distinguishing genuine semantic contributions from
identity proxy effects in such evaluations.

\subsection{Performance Limitations and Baseline Contrasts}

The cross-dataset behavior of FailRank-BB~\cite{hernandes2025method}
is fully explained by the identity proxy effect described above.
On RTPTorrent, FailRank-BB ranks third (0.822) by implicitly learning
per-test failure rates through its embedding fingerprints. On the
industrial dataset, where the test suite is larger, more dynamic, and
severely imbalanced, this implicit learning fails, resulting in the
weakest baseline performance (0.595, $-$27.6\% below \filopriori{}).
The disparity is not attributable to semantic information being more
or less available across datasets; it reflects the effectiveness of
identity-based memorization under different data conditions.

\filopriori{} achieves the highest APFD in 11 of 20 RTPTorrent projects but
underperforms on specific projects relative to the strongest baselines.
Analysis of these cases reveals common characteristics:

\begin{itemize}
    \item \textbf{Sparse failure data}: Projects with very few builds
    containing failures (e.g., 10--11 builds) limit the
    statistical reliability of APFD estimates and reduce the graph's
    effectiveness.
    \item \textbf{High variance}: Projects with inconsistent failure patterns
    exhibit high APFD standard deviations, making reliable prediction
    difficult for any model.
    \item \textbf{Feature sparsity}: In projects with few failures, the
    structural features (failure rate, recency, trend) converge to uninformative
    values, and the co-failure graph is too sparse to provide meaningful signal.
\end{itemize}

In these scenarios, simpler temporal-pattern approaches (TCP-Net, FailRank-BB)
prove more robust for specific projects, as they require less historical data to produce stable
predictions. However, the ensemble
architecture (Section~\ref{sec:ensemble}) effectively mitigates the
worst-case scenarios of standalone graph-based approaches.

To illustrate \filopriori{}'s behavior in practice, we examine two
representative builds from the industrial dataset.
In a high-performance case (Build UUG34.20, APFD = 1.0),
the build contained 4 test cases, all of which failed. \filopriori{} assigned the
highest probability (0.395) to a test case with strong co-failure history
and recent failure patterns, placing all 4 failures at the top of the
ranking, demonstrating the model's ability to exploit
concentrated failure signals through the co-failure graph.
In a low-performance case (Build T2TV33.13, APFD = 0.152),
the build contained 176 test cases with only 4 failures (2.3\%). The model assigned
high probabilities (0.530--0.534) to a cluster of passing tests that shared
strong co-failure and semantic edges with historically unstable tests, while
the actual failures, which were new or had no recent failure history, received
lower scores. This case illustrates a fundamental limitation: when failures
arise from previously unseen patterns (e.g., a new code path or an
infrastructure change), historical features and graph-based propagation
cannot anticipate them.

These cases highlight that \filopriori{}'s effectiveness is contingent on the
presence of recurrent failure patterns. In environments where failures are
predominantly novel (e.g., early development phases or major refactoring
cycles), the graph-based approach provides less advantage over simpler
heuristics.

\subsection{Class Imbalance Handling}

A key design choice of \filopriori{} is the handling of severe class imbalance.
In the industrial dataset, the Pass:Fail ratio is 37:1. The ablation study
reveals that the simplified dual-balancing strategy contributes +4.0\%,
confirming that careful handling of class imbalance is critical for TCP.

The sensitivity analysis shows that the choice of loss function has a 3.6\%
impact on performance. The finding that separating prior correction (balanced
sampling at 29:1) from gradient focusing (focal loss with $\gamma = 2.0$)
outperforms triple compensation ($\approx$646$\times$
effective weight) provides a practical guideline for practitioners: when
multiple imbalance-handling techniques are available, combining
complementary mechanisms operating at different levels (data vs.\ loss)
is preferable to accumulating redundant weights at the same level, which can
cause over-correction and mode collapse.

\subsection{Practical Deployment Considerations}

\filopriori{} can be integrated into CI/CD pipelines to prioritize test
execution. The +10.2\% to +27.6\% improvement over deep learning baselines on
the industrial dataset translates to substantially faster fault detection,
reducing the feedback loop for developers.

In terms of computational cost, training requires approximately 2--3 hours on
a single GPU (NVIDIA RTX 3090), while inference is fast (less than 1 second
per build), making real-time prioritization feasible for most CI environments.
The approach requires historical test execution data with at least 50 builds
for effective graph construction. Projects with limited execution history may
not benefit fully from the graph-based approach, though the DNN ensemble
component can still provide useful prioritization based on per-test temporal
features.

\subsection{Generalization and Automatic Adaptation}
\label{sec:boundary}

A natural question is whether \filopriori{} requires manual configuration for
each new dataset. The framework incorporates two automatic mechanisms that
reduce the need for human intervention.

\textbf{Automatic graph adaptation.} The graph construction methodology
(Section~\ref{sec:graph}) uses only the edge types supported by the available
metadata. Co-failure edges, computed from test execution logs available in any
CI/CD system, form the universal base. When additional metadata exists (commit
messages, component labels, semantic descriptions), the corresponding edge
types are constructed automatically; when it does not, those edge types simply
produce no edges. The practitioner does not need to ``configure'' the graph: the
metadata itself determines the graph structure.

\textbf{Data-driven scoring.} The blending coefficient $\alpha$ between
the GNN and DNN ensemble is optimized automatically on the validation set. On
the industrial dataset, a DNN-only variant achieves APFD = 0.686 ($-$9.9\%
vs.\ the GNN-based model), so the optimizer would naturally favor the GNN
($\alpha \approx 1.0$). On RTPTorrent, the optimizer selects $\alpha \approx 0$
for most projects, naturally favoring the DNN. This data-driven selection
eliminates the need for manual architectural decisions about which scoring
strategy to use.

\textbf{What requires human decision.} In our evaluation, we used two
configurations (Table~\ref{tab:rtptorrent_config}) that differ in training
hyperparameters (learning rate, epochs, batch size) and in whether the KNN
orphan pipeline is active. The hyperparameter differences reflect the different
data scales (277 industrial builds vs.\ up to 42,000 builds per RTPTorrent
project) and could be addressed through automated hyperparameter
search~\cite{bergstra2012random}. The KNN orphan pipeline is an optimization
for the industrial setting where 22.6\% of tests are orphans. In a unified
deployment, the DNN ensemble already covers orphan scoring, making the KNN
pipeline optional rather than essential.

To guide practitioners deploying \filopriori{} on a new dataset, we
summarize the key configuration decisions:
\begin{enumerate}
    \item \textbf{Semantic input}: Determined by data availability: encode
    whatever textual fields exist (test descriptions, test names, commit
    messages). No model change is required; only the input projection
    layer dimension adapts.
    \item \textbf{Graph edge types}: Co-failure edges are computed
    automatically from CI logs; additional edge types (co-success, component,
    semantic, temporal) are constructed when the corresponding metadata exists
    and produce no edges otherwise.
    \item \textbf{DNN ensemble}: We recommend enabling the DNN ensemble by
    default; the validation-optimized $\alpha$ will naturally assign it zero
    weight when the GNN alone is sufficient.
    \item \textbf{Training hyperparameters}: Learning rate and epoch count
    should be scaled to the dataset size (lower learning rate and more epochs
    for smaller datasets). Standard automated tuning (e.g., random
    search~\cite{bergstra2012random}) can replace manual selection.
\end{enumerate}
The remaining architectural parameters (GAT layers, heads, hidden dimensions)
were stable across both evaluation settings and can serve as reasonable
defaults.

\textbf{Boundary conditions and co-failure graph density.} The cross-dataset
results reveal that co-failure graph density is the key factor determining
\filopriori{}'s advantage. On the industrial dataset, where the co-failure
graph is denser (432 edges among 2,283 test cases), \filopriori{} achieves
substantial and statistically significant improvements over all baselines
(+10.2\% to +27.6\%, all $p < 0.001$). On RTPTorrent, where co-failure
graphs are sparser (fewer failure builds per project), the margins narrow
to a statistical tie with the top three baselines after Holm-Bonferroni
correction.

This pattern should be interpreted transparently: the graph-based approach
provides its largest advantage when the co-failure graph has sufficient
density for GATv2 to propagate meaningful failure signals. In repositories
with very sparse failure data (such as some RTPTorrent projects with only
10--11 failure builds), the co-failure graph provides limited structural
information, and simpler temporal-pattern approaches (DeepOrder, TCP-Net)
achieve comparable results. The DNN ensemble component compensates for
this in the adaptive pipeline, becoming the primary contributor when graph
density is low. We recommend that practitioners deploy \filopriori{} with
co-failure edges as the default configuration. The per-edge-type ablation
(Table~\ref{tab:edge_ablation}) confirms this is the optimal setup, and
it requires only CI execution logs, which are universally available.
